\addcontentsline{toc}{section}{Заключение}
\section*{Заключение}


В ходе выполнения лабораторной работы был реализован двумерный клеточный автомат с окрестностью фон Неймана в соответствии с полученным номером 2. Граничный условия - торроидальные. Пользователь определяет ширину поля и количество итераций. Были учтены возможность ввода различных начальных условий (как вручную, так и случайным образом) по выбору пользователя. Был реализован интерфейс. Также клеточный автомат проанализирован, в отчете зафиксированно его поведение, паттерны, "сходимость" и энтропия.

\subsection*{Достоинства}

Приложение позволяет менять состояние клеток вручную, выбирать случайное
начальное заполнение и наблюдать выполнение одного или нескольких шагов.
Счётчики живых и мёртвых клеток упрощают наблюдение за изменением поля.

\subsection*{Недостатки}

Начальное случайное состояние нельзя сохранить или повторить из интерфейса,
так как значение генератора случайных чисел не задаётся пользователем.
Также в программе нет сохранения поля и экспорта результатов.

\subsection*{Масштабирование}

Для поля размером $n\times n$ временная сложность одного шага и сложность
по памяти составляют $O(n^2)$. При увеличении стороны поля в два раза число
клеток и объём вычислений возрастают примерно в четыре раза. Для работы с
более крупными полями потребуется оптимизировать обновление и отрисовку.
